% =====================================================================
%  A 题 药材的烘干问题 —— 模型公式（LaTeX 源码）
%  编译： xelatex formulas.tex        （需要 TeX Live 的 ctex 宏包）
%  或：   pdflatex formulas.tex       （若 ctex 不可用，见文末 ctex 备选）
% =====================================================================
\documentclass[11pt]{ctexart}
\usepackage{amsmath,amssymb,bm}
\usepackage{geometry}
\usepackage{booktabs}
\geometry{a4paper,margin=2.4cm}
\allowdisplaybreaks

\newcommand{\dd}{\mathrm{d}}
\newcommand{\pd}[2]{\frac{\partial #1}{\partial #2}}

\title{A 题\quad 药材的烘干问题\\[4pt]\large 数学模型与经验公式汇总}
\author{}
\date{}

\begin{document}
\maketitle

\section{控制方程（柱坐标，径向一维）}

药材为半径 $R_0=2\ \mathrm{cm}$、长 $L=25\ \mathrm{cm}$ 的圆柱体。取 $r$ 为到药材中心的距离，
只保留径向变化。热量守恒（Fourier 定律）与水分守恒（Fick 定律，$X$ 为干基含水率）：

\begin{equation}
\rho c_p \pd{T}{t}
=\frac{1}{r}\pd{}{r}\!\left(k\,r\pd{T}{r}\right),
\qquad 0<r<R(t),\ t>0,
\end{equation}

\begin{equation}
\pd{X}{t}
=\frac{1}{r}\pd{}{r}\!\left(D\,r\pd{X}{r}\right),
\qquad 0<r<R(t),\ t>0 .
\end{equation}

当 $R(t)\equiv R_0$ 时（问题 1--3），上两式即固定区域上的轴对称抛物方程。

\section{定解条件}

\subsection*{中心对称条件（$r=0$）}
\begin{equation}
\left.\pd{T}{r}\right|_{r=0}=0,
\qquad
\left.\pd{X}{r}\right|_{r=0}=0 .
\end{equation}

\subsection*{表面第三类边界（$r=R(t)$）}
\begin{equation}
-k\left.\pd{T}{r}\right|_{r=R}=h\bigl(T(R,t)-T_f(t)\bigr),
\end{equation}
\begin{equation}
-D\left.\pd{X}{r}\right|_{r=R}=h_m\bigl(X(R,t)-X_f(t)\bigr).
\end{equation}
其中 $T_f,X_f$ 为烘房（外部空气）的温度与水分浓度，$h$ 为对流换热系数，
$h_m$ 为对流传质系数。量纲检验：式 (5) 左端为 $\mathrm{(m^2/s)\cdot(kg/kg)/m}$，
右端为 $\mathrm{(m/s)\cdot(kg/kg)}$，一致。

\subsection*{初始条件}
\begin{equation}
T(r,0)=T_0=28\ ^\circ\mathrm{C},\qquad
X(r,0)=X_0=2.55\ \mathrm{kg/kg},\qquad 0\le r\le R_0 .
\end{equation}

\section{物性与经验公式}

\subsection*{附录 2（问题 1）}
\begin{equation}
\rho=820\ \tfrac{\mathrm{kg}}{\mathrm m^3},\quad
c_p=2600\ \tfrac{\mathrm J}{\mathrm{kg\cdot K}},\quad
k=0.36\ \tfrac{\mathrm W}{\mathrm{m\cdot K}},\quad
h=25\ \tfrac{\mathrm W}{\mathrm{m^2\cdot K}},\quad
h_m=8\times10^{-7}\ \tfrac{\mathrm m}{\mathrm s},
\end{equation}
\begin{equation}
\boxed{\;D=7\times10^{-9}\exp\!\left(-\frac{0.89}{X}\right)\ \ \mathrm{m^2/s}\;}
\end{equation}

\subsection*{附录 3（问题 2、问题 3）}
\begin{equation}
D=2.4\times10^{-3}\exp\!\left(-\frac{0.45}{X}\right)\exp\!\left(-\frac{3850}{T_K}\right)
\ \ \mathrm{m^2/s},
\end{equation}
\begin{equation}
\rho=650+128X\ \ \mathrm{kg/m^3},
\qquad
c_p=1450+2736\,\frac{X}{X+1}\ \ \mathrm{J/(kg\cdot K)},
\qquad
k=0.21+0.38\,\frac{X}{X+1}\ \ \mathrm{W/(m\cdot K)} .
\end{equation}

\subsection*{附录 4（问题 4）}
\begin{equation}
D=4.2\times10^{-4}\exp\!\left(-\frac{0.30}{X}\right)\exp\!\left(-\frac{3850}{T_K}\right)
\ \ \mathrm{m^2/s},
\end{equation}
\begin{equation}
\rho=760+90X\ \ \mathrm{kg/m^3},
\qquad
c_p=1850+2150\,\frac{X}{X+1}\ \ \mathrm{J/(kg\cdot K)},
\qquad
k=0.12+0.20\,\frac{X}{X+1}\ \ \mathrm{W/(m\cdot K)} .
\end{equation}
其中 $T_K=T+273.15$ 为热力学温度（单位 K）。

\section{环境条件（两段式）}
由附件 1 的实测数据线性插值得到预热平衡段，其后为恒温干燥段：
\begin{equation}
T_f(t)=
\begin{cases}
\text{附件 1 线性插值}, & 0\le t\le t_1=14400\ \mathrm s,\\[2pt]
50\ ^\circ\mathrm C, & t>t_1,
\end{cases}
\qquad
X_f(t)=
\begin{cases}
\text{附件 1 线性插值}, & 0\le t\le t_1,\\[2pt]
0.05\ \mathrm{kg/kg}, & t>t_1 .
\end{cases}
\end{equation}

\section{问题 4 的 Landau（前沿固定）变换}

令 $\xi=r/R(t)\in[0,1]$、$\tau=t$。链式法则给出
\begin{equation}
\pd{}{r}=\frac1R\pd{}{\xi},
\qquad
\left.\pd{}{t}\right|_r=\left.\pd{}{\tau}\right|_{\xi}-\frac{\xi \dot R}{R}\pd{}{\xi},
\qquad
\frac1r\pd{}{r}\!\left(rD\pd{}{r}\right)
=\frac{1}{R^2\xi}\pd{}{\xi}\!\left(\xi D\pd{}{\xi}\right).
\end{equation}

空间坐标下含固相牵连速度 $v_s$ 的干基守恒方程为
\begin{equation}
\pd{X}{t}+v_s\pd{X}{r}
=\frac1r\pd{}{r}\!\left(rD\pd{X}{r}\right).
\end{equation}
变换后
\begin{equation}
\pd{X}{\tau}+\frac{v_s-\xi\dot R}{R}\pd{X}{\xi}
=\frac{1}{R^2\xi}\pd{}{\xi}\!\left(\xi D\pd{X}{\xi}\right).
\end{equation}
仿射（均匀）收缩假定下 $r(t)=\xi R(t)$，故
\begin{equation}
v_s=\left.\frac{\dd r}{\dd t}\right|_{\xi}=\xi\dot R
\quad\Longrightarrow\quad
\frac{v_s-\xi\dot R}{R}\equiv 0 ,
\end{equation}
对流项恒等于零，方程化为
\begin{equation}
\boxed{\;
\pd{X}{\tau}=\frac{1}{R^2\xi}\pd{}{\xi}\!\left(\xi D\pd{X}{\xi}\right),
\qquad
\rho c_p\pd{T}{\tau}=\frac{1}{R^2\xi}\pd{}{\xi}\!\left(\xi k\pd{T}{\xi}\right)
\;}
\end{equation}
即与固定区域方程同形，只是 $R=R(\tau)$ 随时间变化。

表面条件在 $\xi=1$ 处（注意 $\partial/\partial r=R^{-1}\partial/\partial\xi$）：
\begin{equation}
k\left.\pd{T}{\xi}\right|_{\xi=1}=Rh\bigl(T_f-T\bigr),
\qquad
D\left.\pd{X}{\xi}\right|_{\xi=1}=Rh_m\bigl(X_f-X\bigr).
\end{equation}

半径 $R(t)$ 由附件 2 用保形分段三次 Hermite 插值（PCHIP）给出：
\begin{equation}
R(t)=\mathrm{PCHIP}\{t_i,R_i\}(t),
\qquad
\dot R(t)=\frac{\dd}{\dd t}\mathrm{PCHIP}\{t_i,R_i\}(t).
\end{equation}

\section{数值离散}

\subsection*{有限体积离散（乘以 $\xi$ 后积分）}
内部节点 $j=1,\dots,N-1$：
\begin{equation}
\frac{\dd X_j}{\dd \tau}
=\frac{1}{R^2\xi_j\Delta\xi^2}
\Bigl[\xi_{j+\frac12}D_{j+\frac12}\bigl(X_{j+1}-X_j\bigr)
-\xi_{j-\frac12}D_{j-\frac12}\bigl(X_j-X_{j-1}\bigr)\Bigr].
\end{equation}
中心节点（用 L'H\^opital 极限 $\xi^{-1}\partial_\xi(\xi\partial_\xi X)\to2\partial_\xi^2X$）：
\begin{equation}
\frac{\dd X_0}{\dd\tau}=\frac{4D_{1/2}}{R^2\Delta\xi^2}\bigl(X_1-X_0\bigr).
\end{equation}
表面节点（控制体 $[\xi_{N-\frac12},1]$，权重 $w=1-\xi_{N-\frac12}^2$）：
\begin{equation}
\frac{\dd X_N}{\dd\tau}
=\frac{2}{R^2w}\left[Rh_m\bigl(X_f-X_N\bigr)
-\frac{\xi_{N-\frac12}D_{N-\frac12}}{\Delta\xi}\bigl(X_N-X_{N-1}\bigr)\right].
\end{equation}
界面扩散系数取调和平均：
\begin{equation}
D_{j+\frac12}=\frac{2D_jD_{j+1}}{D_j+D_{j+1}} .
\end{equation}

\subsection*{时间离散}
扩散项用后向欧拉（全隐式，无条件稳定），第 $n+1$ 时间层：
\begin{equation}
\frac{X_j^{n+1}-X_j^{n}}{\Delta\tau}
=\frac{1}{R^2\xi_j\Delta\xi^2}
\Bigl[\xi_{j+\frac12}\tilde D_{j+\frac12}\bigl(X_{j+1}^{n+1}-X_j^{n+1}\bigr)
-\xi_{j-\frac12}\tilde D_{j-\frac12}\bigl(X_j^{n+1}-X_{j-1}^{n+1}\bigr)\Bigr],
\end{equation}
其中 $\tilde D$ 为 Picard 迭代的当前值。整理为三对角方程组
$a_jX_{j-1}+b_jX_j+c_jX_{j+1}=d_j$，用 Thomas 追赶法 $O(n)$ 求解。

\section{关键结果}

\begin{table}[h]\centering
\begin{tabular}{ll}
\toprule
问题 & 结果 \\
\midrule
问题 1（1800 s） & 温度 $33.5765/33.7732/34.3653/35.3631/36.7863\ ^\circ$C\\
                 & 含水率 $2.5500/2.5497/2.5383/2.3755/1.5103$ kg/kg\\
问题 2（3 h）    & 温度 $49.8494/49.8553/49.8745/49.9101/49.9664\ ^\circ$C\\
                 & 含水率 $1.7662/1.7166/1.5702/1.3333/1.0081$ kg/kg\\
问题 3           & $t_{\text{dry}}=206958\ \mathrm s=57.49\ \mathrm h=2.395$ 天\\
问题 4           & $t_{\text{dry}}=183926\ \mathrm s=51.09\ \mathrm h=2.129$ 天，$R_{\text{末}}=1.200$ cm\\
\bottomrule
\end{tabular}
\end{table}

\end{document}
